% ============================================================
% ［架空の記入例・提出不可］研究上の位置づけと適用条件は説明用である．
% 引用する4件の文献は実在する．
% 通常のchapters/02-background.texには自動で読み込まれない．
% ============================================================
\chapter{背景・関連研究}

\begin{center}
  \UsamiFictionalExampleNotice
\end{center}

\section{腹部CT多臓器分割}

医用画像の領域分割では，画像の各画素またはボクセルを解剖構造のクラスへ割り当てる．U-Net\cite{ronneberger2015unet}は，縮小経路で抽出した文脈情報と拡大経路の位置情報をskip connectionで結合する．3D U-Net\cite{cicek20163dunet}は3次元畳み込みを用い，スライス間の連続性を含む体積情報からボクセル単位の予測を行う．腹部CTでは，肝臓のような大きな臓器と，膵臓や細い血管のような小さく境界が不明瞭な構造を同時に扱うため，クラス不均衡と境界位置の不確実性が生じる．

分割器の評価には，予測領域と正解領域の重複を測るDice係数が広く用いられる．ただし，少数の離れた誤検出があっても体積の大きな臓器ではDiceへの影響が小さい場合がある．そこで本例では，境界間距離の上側分位点を表すHD95も併記し，領域重複と局所的な境界逸脱を別の軸で評価する．

\section{拡散モデルによるデータ拡張}

Denoising Diffusion Probabilistic Models\cite{ho2020ddpm}は，データへ段階的に雑音を加える前向き過程と，雑音からデータを復元する逆過程を学習する生成モデルである．高解像度の3次元体積を画像空間で直接生成すると計算量が大きいため，本例では潜在表現上で拡散過程を学習する潜在拡散モデル\cite{rombach2022ldm}の考え方を用いる．臓器ラベルを条件として与えることで，指定した構造に対応するCT体積を生成する架空の構成を想定する．

分割学習用のデータ拡張では，生成画像の写実性だけでなく，条件ラベルと画像中の臓器境界が対応している必要がある．生成された臓器の位置や体積が学習分布から大きく外れる場合，合成対をそのまま加えると分割器が解剖学的に不自然な対応を学習する可能性がある．したがって，生成品質と教師ラベルの妥当性を分けて検査し，不適切な合成対を学習前に除く設計が必要となる．

\section{解剖構造の保持}

腹部臓器には，左右の腎臓の位置関係，臓器ごとの体積範囲，互いに重複しない領域等の制約がある．本例では，これらを厳密な解剖モデルとして固定せず，学習集合から得た分布と明示的な非重複条件を合成対の品質管理に用いる．さらに，分割器の予測に対しても，隣接関係からの過度な逸脱を抑制する正則化項を加える．生成時の採択と分割時の正則化は目的が異なるため，除去実験によって個別の寄与を確認する．

\section{本研究の位置づけ}

本研究は，U-Net\cite{ronneberger2015unet}および3D U-Net\cite{cicek20163dunet}の分割構造そのものを新規提案するものではない．また，拡散モデル\cite{ho2020ddpm}や潜在拡散モデル\cite{rombach2022ldm}の一般的な生成能力を比較する研究でもない．少数症例の腹部CT分割に用途を限定し，合成画像と条件ラベルの解剖学的対応を検査した上で，採択済み合成対を分割学習へ接続する点に焦点を置く．本章の適用条件と差分は架空の記入例であり，既存研究に対する実在の新規性主張ではない．
